\origpage{116}{ROBERTSON AND SEYMOUR}
\hspace*{12pt}(ii) \textit{there is a $\Sigma'$-span of order $\phi$, where $\Sigma'$ is a surface obtained by adding a crosscap to $\Sigma$, or}

\hspace*{12pt}(iii) \textit{there is a $\mathscr{T}^*$-central segregation of $G$ of type $(\rho_k, k)$ with an arrangement in $\Sigma$.}

\textit{Proof of} (1.1), \textit{assuming} (1.2). Let $\Sigma, \kappa, \phi, \mu$ be as in (1.1). We may assume that $\mu \ge 12$ and is even, for every $(\lambda, \mu)$-flat $\Sigma$-span is also $(\lambda, \mu')$-flat for all $\mu' \ge \mu$ (see (5.1)). We may also assume that $\phi \ge 8$, for if (1.1) holds for $\phi \ge 8$ then it holds for all $\phi$. Define $\theta_{\delm{k}\fix{\kappa}+1} = \max(\phi, \mu+16)$ and $\lambda_{\delm{k}\fix{\kappa}+1} = 4$;\fixnote{fix:1.1-proof-theta}{No $k$ is bound here, and the induction that follows starts from $\theta_{\kappa+1}$ and $\lambda_{\kappa+1}$.} and for $k = \kappa, \kappa-1, \dots, 0$, define $\theta_k, \lambda_k, \rho_k$ inductively by (1.2). Let $\theta = \theta_0$, $\lambda = \lambda_0$ and $\rho = \max(\rho_0, \dots, \rho_\kappa)$. We claim that $\theta, \lambda, \rho$ satisfy (1.1). For let $G, \mathscr{T}^*$ be as in (1.1); then there is a $(\lambda_0, \mu)$-flat $\Sigma$-span of order $\ge \theta_0$. Choose $k \ge 0$ with $k \le \kappa+1$ maximum such that there is a $(\lambda_k, \mu)$-flat $\Sigma$-span of order $\ge \theta_k$ (with order $\theta_k$ if $k \ge 1$), with $k$ eyes mutually at distance $\ge \theta_k$. Suppose first that $k \le \kappa$. Since (1.2)(i) does not hold, it follows from (1.2) that either (1.2)(ii) or (1.2)(iii) holds, and hence either (1.1)(ii) or (1.1)(iii) holds (because a segregation of type $(\rho_k, \kappa)$ is also of type $(\rho, \kappa)$). We may therefore assume that $k = \kappa+1$. Let $H, \eta, \mathscr{T}$ be a $(\lambda_{\delm{k}\fix{\kappa}+1}, \mu)$-flat $\Sigma$-span of order $\theta_{\kappa+1}$,\fixnote{fix:1.1-proof-lambda}{With $k = \kappa+1$, $\lambda_{k+1}$ would be $\lambda_{\kappa+2}$, but the span given by the choice of $k$ is $(\lambda_k, \mu)$-flat, that is, $(\lambda_{\kappa+1}, \mu)$-flat, and the next sentence uses $\lambda_{\kappa+1} = 4$.} with $\kappa+1$ eyes $r_1, \dots, r_{\kappa+1}$ mutually at distance $\theta_{\kappa+1}$. Since $H, \eta, \mathscr{T}$ is $(4, \mu)$-flat, and hence $(0, \mu)$-flat (see (5.1)), there is no $\eta(H)$-path in $G$ with ends $\eta(u), \eta(v)$, where $u \in V(H) \cap \bar{r}_i$, $v \in V(H) \cap \bar{r}_j$, for $i \ne j$, since $d(u, v) \ge \mu$ because
\begin{equation*}
\mu+2 \le \theta_{\kappa+1} = d(r_i, r_j) \le d(u, r_i) + d(u, v) + d(v, r_j) \le d(u, v) + 2.
\end{equation*}
Hence $r_1, \dots, r_{\kappa+1}$ are independent eyes, and so (1.1)(i) holds. $\square$

\sectionhead{2. \textsc{Tangles} \textsc{And} \textsc{Transactions}}

To prove (1.2) we need at some stage to prove that a certain segregation is $\mathscr{T}^*$-central, and to do so we shall use the following lemma.

\stmt
(2.1) \textit{Let $\mathscr{T}$ be a tangle in a graph $G$, let $\rho \ge 1$ with $ord(\mathscr{T}) \ge 5\rho+2$, and let $\mathscr{S}$ be a segregation of $G$ of type $(\rho, \kappa)$ for some $\kappa$. Suppose that for every $(C, \Omega) \in \mathscr{S}$ there is no $(A, B) \in \mathscr{T}$ of order $\le 2\rho+1$ with $B \subseteq C$. Then $\mathscr{S}$ is $\mathscr{T}$-central.}

\textit{Proof.} Let $(C, \Omega) \in \mathscr{S}$, and suppose, for a contradiction, that $(A, B) \in \mathscr{T}$ and $B \subseteq C$. If $|\Omega| \ge 4$ then $(C, \Omega)$ is a $\rho$-vortex since $\mathscr{S}$ is of type $(\rho, \kappa)$. If $|\bar{\Omega}| \le 3$ then $(C, \Omega)$ has no transaction of cardinality $\ge 2$, and again it is a $\rho$-vortex since $\rho \ge 1$. \fix{We may assume that $\bar{\Omega} \neq \emptyset$.}\fixnote{fix:2.1-proof-empty}{A gap in the proof: if $\bar{\Omega} = \emptyset$ and $C$ is not null, Graph Minors IX, (8.1) does not apply, since a linear decomposition needs $n = |\bar{\Omega}| \ge 1$ (see the note to (8.1) in Graph Minors IX). In that case $C$ shares no vertex with the rest of $G$, so there is a separation of $G$ of order $0 \le 2\rho + 1$ whose second part is $C$. The hypothesis of (2.1), with the tangle axioms, then rules out $(A, B) \in \mathscr{T}$ with $B \subseteq C$ directly. So (2.1) holds as printed.} Hence in either case, from [1, Theorem (8.1)] we may enumerate $\bar{\Omega} = \{t_1, \dots, t_n\}$ in order so that there are subgraphs $X_1, \dots, X_n$ of $C$ satisfying (1)--(4) below:
